% !TeX root = ../常微分方程笔记.tex
% 由主文档通过 \input 载入；请编译主文档。
\setcounter{chapter}{0} % 保留原书章号。
\chapter{绪论与基本概念}
\label{chap:ode-intro}

微分方程用局部变化规律描述整体演化。学习时需要同时回答三个问题：模型表达了什么，方程有哪些解，这些解具有怎样的性质。显式公式是重要工具，但并不是研究的唯一目标。

\section{从变化规律到微分方程}
设 $N(t)$ 表示时刻 $t$ 的种群数量。在资源充足、单位个体净增长率保持为常数 $r$ 的假设下，短时间内的净增长量近似为 $rN(t)\Delta t$；令 $\Delta t\to0$，得到
\[
  N'=rN,\qquad N(0)=N_0.
\]
其中 $r$ 的单位是时间的倒数。由 $(e^{-rt}N)'=0$ 得 $N(t)=N_0e^{rt}$。指数增长是模型假设的结果；若资源限制改变了单位增长率，方程也必须改变。

类似地，物体温度 $T(t)$ 满足牛顿冷却模型
\[
 T'=-k(T-T_{\mathrm{env}}),\qquad k>0,
\]
其中环境温度 $T_{\mathrm{env}}$ 暂取常数。弹簧振子的位移 $y(t)$ 则满足
\[
 my''+cy'+ky=F(t),\qquad m>0,\ c\geq0,\ k>0.
\]
此处 $m,c,k$ 分别为质量、阻尼系数和弹性系数，$F$ 为外力；$-ky$ 与 $-cy'$ 分别是线性回复力和阻尼力。方程中的每一项均具有力的量纲。模型是否适用取决于这些力学近似，而不取决于能否解出方程。

\begin{example}[冷却模型]
一物体初温为 $80\,{}^\circ\mathrm C$，置于恒温 $20\,{}^\circ\mathrm C$ 的环境中。以分钟为时间单位，若十分钟后温度为 $50\,{}^\circ\mathrm C$，求模型中的 $k$ 及后续温度。
\end{example}
\begin{solve}
令 $u=T-20$，则 $u'=-ku$，所以 $T(t)=20+60e^{-kt}$。代入 $T(10)=50$ 得 $e^{-10k}=1/2$，从而
\[
 k=\frac{\log2}{10},\qquad T(t)=20+60\cdot2^{-t/10}.
\]
温差每十分钟减半；在该模型中温度渐近接近环境温度，而不在有限时间内与之相等。
\end{solve}

\section{方程的阶、线性与自治性}
\begin{definition}
\emph{常微分方程}（ordinary differential equation，ODE）是未知函数及其关于同一个自变量的导数之间的关系。一般标量形式为
\[
 F(t,y,y',\ldots,y^{(n)})=0.
\]
实际出现的最高导数阶数 $n$ 称为方程的阶。若可解成 $y^{(n)}=g(t,y,\ldots,y^{(n-1)})$，称为关于最高阶导数的标准形式。未知函数依赖多个自变量且涉及偏导数时，得到偏微分方程（partial differential equation，PDE）。
\end{definition}

方程
\begin{equation}
 a_n(t)y^{(n)}+a_{n-1}(t)y^{(n-1)}+\cdots+a_0(t)y=b(t)
 \label{eq:ode-general-linear}
\end{equation}
称为线性方程（linear equation），因为未知函数及其导数以线性组合出现。$b\equiv0$ 时称齐次，否则称非齐次。讨论标准形时应限定在 $a_n\ne0$ 的区间。$y''+ty=\sin t$ 是线性的，$y'=y^2$ 和 $y''+yy'=0$ 不是线性的。

标准形右端不显含 $t$ 时称为自治方程（autonomous equation）。例如 $y'=y^2$ 自治而非线性，$y'=ty$ 线性而非自治。自治性与线性是不同分类标准。一阶方程 $y'=G(y/t)$ 的“齐次型”是另一术语，不能与线性齐次方程混淆。

\section{解、解区间与隐式表示}
\begin{definition}
设 $I$ 为非空开区间。若 $y\in C^n(I)$，且将 $y(t)$ 及其导数代入方程后对所有 $t\in I$ 均成立，则称 $y$ 是方程在 $I$ 上的经典解。解应连同其定义区间一起给出。
\end{definition}

显式式子 $y=\varphi(t)$ 直接给出函数；隐式关系 $H(t,y)=0$ 需要选定可微分支。当 $H\in C^1$ 且某点处 $H_y\ne0$ 时，隐函数定理保证该点附近可以解出 $y$。例如 $y^2=t$ 在 $t>0$ 有两条分支，而它们在 $t=0$ 处均没有有限右导数。

由积分常数组成并覆盖所讨论解族的表达式称为通解，其中固定常数得到特解。对于正则 $n$ 阶线性方程，通解含 $n$ 个独立常数；一般隐式非线性方程可能有多分支或不包含在通常参数族中的奇解，不能只用“常数个数”定义通解。

\begin{example}[表达式与最大区间]
求 $y'=y^2$、$y(0)=1$ 的解及其包含零点的最大存在区间。
\end{example}
\begin{solve}
在 $y\ne0$ 时，$(-1/y)'=1$，故 $y=1/(C-t)$。初值给出 $C=1$，因而 $y(t)=1/(1-t)$。它在 $(-\infty,1)$ 上满足方程与初值，并在 $t\uparrow1$ 时趋于正无穷，故不能以有限连续值延伸过 $1$。相同表达式在 $(1,\infty)$ 上也满足方程，但那是另一段解，不是给定初值解的一部分。原方程另有常值解 $y\equiv0$。
\end{solve}

\section{初值问题与边值问题}
\emph{初值问题}（initial value problem，IVP）在同一时刻指定状态。例如
\begin{equation}
 y'=f(t,y),\qquad y(t_0)=y_0
 \label{eq:ode-scalar-ivp}
\end{equation}
的一般任务是确定经过 $(t_0,y_0)$ 的解。$n$ 阶标准方程通常给定 $y(t_0),y'(t_0),\ldots,y^{(n-1)}(t_0)$。\emph{边值问题}（boundary value problem，BVP）则在不同端点施加条件，例如 $y''+y=0$、$y(0)=a$、$y(\pi)=b$。

即使初值问题唯一可解，边值问题也可能无解或有无穷多解。上述方程的通解是 $A\cos t+B\sin t$；边界条件要求 $A=a$、$-A=b$。若 $b\ne-a$ 则无解，若 $b=-a$ 则 $B$ 任意。存在性、唯一性和数据扰动下的连续依赖性都需要各自的条件。

\section{方向场与解的几何意义}
一阶方程 $y'=f(t,y)$ 在平面每一点指定斜率 $f(t,y)$，这些短线元组成方向场（slope field）。解的图像是与方向场相切的积分曲线（integral curve）。在 $f>0$ 的区域解增加，在 $f<0$ 的区域解减少；$f=0$ 只表示该点切线水平，不一定意味着整条解为常数。

对于 $y'=g(y)$，满足 $g(y_*)=0$ 的常值 $y_*$ 是平衡点（equilibrium）。例如 $y'=y(1-y)$ 的平衡点为 $0,1$；在 $0<y<1$ 上解增加，在 $y>1$ 上解减少。若局部唯一性成立，非平衡解不能在有限时间穿过平衡解；这一结论将在第\ref{chap:ode-existence}章得到理论依据。

\section{习题}
\begin{enumerate}
 \item\label{ex:ode-classify} 判别 $y''+t^2y=0$、$y'=\sin y$、$y'=t+y^2$ 的阶数、线性与自治性。
 \item\label{ex:ode-square} 求 $y'=y^2$、$y(0)=a$ 的解及最大区间，分别讨论 $a>0$、$a=0$、$a<0$。
 \item\label{ex:ode-boundary-intro} 对 $y''=0$ 比较条件 $y(0)=1,y'(0)=2$ 与 $y(0)=1,y(1)=2$ 所确定的解。
 \item 说明“右端在某点为零”为什么不必使经过该点的解恒为常数，并给出非自治方程的例子。
\end{enumerate}

